\documentclass[11pt,a4paper]{article}
\usepackage[utf8]{inputenc}
\usepackage[T1]{fontenc}
\usepackage{amsfonts}
\usepackage[french]{babel}
\frenchbsetup{StandardLists=true}
\usepackage{enumitem}
\usepackage{amssymb}
\usepackage{amsmath}
\usepackage[left=2cm,right=2cm,top=2cm,bottom=2cm]{geometry}
\usepackage{multirow}
\usepackage[dvipsnames]{xcolor}
\usepackage{fouriernc}
\usepackage{graphicx}
\usepackage{setspace}
\singlespacing
\usepackage{tcolorbox}
\tcbuselibrary{skins, breakable}
\usepackage{multicol}
\usepackage{caption}
\usepackage{fancyhdr}
\pagestyle{fancy}
\renewcommand\headrulewidth{1pt}
\fancyhead[L]{Lycée Denis Diderot}
\fancyhead[R]{Classe de 3\textsuperscript{ème}}
\setcounter{page}{1}\fancyfoot[C]{page \thepage}
\fancyfoot[R]{Année 2025-2026}
\fancyfoot[L]{Alexis RUIZ}
\usepackage{multirow,tabularx,booktabs}
\newcommand{\cadreblanc}[1]{%
\medskip\framebox[\linewidth][c]{%
\rule{0mm}{#1mm}}\par
\medskip}
\usepackage{awesomebox}
\setlength{\aweboxvskip}{0mm}
\usepackage{pifont}
\setlength{\parindent}{0pt}
\setlength{\headheight}{14pt}
\frenchbsetup{StandardLists=true}
\renewcommand{\arraystretch}{1.8}

% Boîte pour les documents documentaires
\newtcolorbox{docdoc}[2]{
  colback=blue!5!white,
  colframe=blue!55!black,
  title={\textbf{Document #1} \quad -- \quad \textit{#2}},
  fonttitle=\bfseries\normalsize,
  breakable,
  before upper={\setlength{\parindent}{0pt}},
}

\begin{document}

%------------------------------------------------------------
% TITRE
%------------------------------------------------------------
\begin{center}
  {\LARGE\textbf{Activité documentaire --- Introduction à l'électricité}}
\end{center}

\hrule
\vspace{3mm}
\noindent Nom : \dotfill \quad Prénom : \dotfill \quad Date : \dotfill
\vspace{3mm}
\hrule
\vspace{5mm}

\begin{center}
  \textit{L'électricité est omniprésente dans notre vie quotidienne. À travers ces trois documents, tu vas découvrir d'où vient l'électricité, comment fonctionne un circuit électrique, et quels dangers il faut éviter.}
\end{center}

\vspace{5mm}

%------------------------------------------------------------
% DOCUMENT 1 : L'électricité dans notre quotidien
%------------------------------------------------------------

\begin{docdoc}{1}{L'électricité dans notre quotidien}

L'électricité est une forme d'énergie indispensable à notre société moderne. En France, chaque habitant consomme en moyenne \textbf{7\,000 kWh} d'électricité par an. Elle est utilisée dans quasiment tous les domaines :

\begin{itemize}
  \item \textbf{L'éclairage} : ampoules LED, lampadaires, néons dans les bureaux.
  \item \textbf{L'électroménager} : réfrigérateur, four, lave-linge, aspirateur.
  \item \textbf{Les transports} : voitures électriques, tramways, trains (TGV).
  \item \textbf{La communication} : téléphones, ordinateurs, internet, télévision.
  \item \textbf{La médecine} : IRM, scanners, appareils de bloc opératoire.
\end{itemize}

En France, l'électricité est produite principalement par des \textbf{centrales nucléaires} (environ 70\,\%), mais aussi par des éoliennes, des panneaux solaires et des barrages hydrauliques. Elle est transportée sur le territoire par des lignes à \textbf{haute tension} (jusqu'à 400\,000\,V), puis transformée pour être distribuée dans les foyers à \textbf{230\,V}.

\vspace{2mm}
\footnotesize\textit{Source : RTE (Réseau de Transport d'Électricité) --- bilan électrique 2024.}
\normalsize
\end{docdoc}

\vspace{4mm}

%------------------------------------------------------------
% DOCUMENT 2 : Les composants d'un circuit électrique
%------------------------------------------------------------

\begin{docdoc}{2}{Les composants d'un circuit électrique}

Pour qu'un circuit électrique fonctionne, quatre types de composants sont nécessaires :

\begin{center}
\renewcommand{\arraystretch}{1.6}
\begin{tabularx}{\linewidth}{|l|X|X|}
\hline
\textbf{Composant} & \textbf{Rôle} & \textbf{Exemples} \\
\hline
Générateur & Fournit l'énergie électrique & Pile, batterie, panneau solaire \\
\hline
Récepteur & Convertit l'énergie électrique en une autre forme & Lampe (lumière), moteur (mouvement) \\
\hline
Interrupteur & Ouvre ou ferme le circuit & Bouton ON/OFF, interrupteur mural \\
\hline
Fils conducteurs & Permettent la circulation du courant & Fils de cuivre, câbles \\
\hline
\end{tabularx}
\end{center}

\vspace{3mm}

Un circuit électrique est dit \textbf{fermé} lorsque le courant peut circuler en continu. Si le circuit est \textbf{ouvert} (interrupteur ouvert, fil coupé), le courant ne circule plus et le récepteur ne fonctionne pas.

\vspace{2mm}
Le courant électrique est un \textbf{déplacement ordonné de charges électriques} dans un conducteur. Par convention, le courant circule du pôle \textbf{positif (+)} vers le pôle \textbf{négatif ($-$)} du générateur, à l'extérieur de celui-ci.

\end{docdoc}

\vspace{4mm}

%------------------------------------------------------------
% DOCUMENT 3 : Dangers de l'électricité et sécurité
%------------------------------------------------------------

\begin{docdoc}{3}{Les dangers de l'électricité et les règles de sécurité}

L'électricité est indispensable mais peut être \textbf{dangereuse} si l'on ne respecte pas certaines règles.

\textbf{Les dangers :}
\begin{itemize}
  \item L'\textbf{électrocution} : le passage du courant dans le corps peut être mortel. À partir de \textbf{25\,V en courant alternatif} (réseau domestique), la tension devient dangereuse pour l'être humain. Le réseau électrique domestique est à 230\,V : il est donc très dangereux.
  \item L'\textbf{incendie} : un court-circuit ou une surcharge peut enflammer les fils et provoquer un incendie.
  \item Les \textbf{brûlures} : un arc électrique (étincelle) peut provoquer des brûlures graves.
\end{itemize}

\textbf{Règles de sécurité à respecter :}
\begin{itemize}
  \item Ne jamais toucher une prise électrique avec les mains mouillées.
  \item Ne jamais introduire d'objets métalliques dans une prise.
  \item Toujours \textbf{couper le courant} (disjoncteur) avant d'intervenir sur une installation.
  \item Utiliser des appareils aux normes (logos \textbf{NF} ou \textbf{CE}).
  \item En cas d'accident, ne jamais toucher la victime directement : couper le courant d'abord, puis appeler le \textbf{15} (SAMU) ou le \textbf{18} (pompiers).
\end{itemize}

\footnotesize\textit{Source : INRS (Institut National de Recherche et Sécurité) --- guide sécurité électrique, 2024.}
\normalsize

\end{docdoc}

\vspace{5mm}

%------------------------------------------------------------
% QUESTIONS
%------------------------------------------------------------
\hrule
\vspace{3mm}
{\Large\textbf{Questions}}
\vspace{3mm}
\hrule
\vspace{5mm}

\begin{enumerate}[label=\textbf{Q\arabic*.}, leftmargin=2cm, itemsep=6pt]

  %--- Document 1 ---
  \item \textit{(Document 1)} --- Cite \textbf{trois domaines} dans lesquels l'électricité est utilisée dans notre quotidien.

  \cadreblanc{25}

  \item \textit{(Document 1)} --- Quelle est la principale source de production d'électricité en France ? Quel pourcentage représente-t-elle ?

  \cadreblanc{22}

  \newpage
  \item \textit{(Document 1)} --- À quelle tension l'électricité est-elle distribuée dans les foyers français ? Comment est-elle transportée avant d'arriver chez nous ?

  \cadreblanc{22}

  %--- Document 2 ---
  \item \textit{(Document 2)} --- Quel est le rôle d'un générateur dans un circuit électrique ? Donne deux exemples de générateurs.

  \cadreblanc{22}

  \item \textit{(Document 2)} --- Complète le tableau suivant :

  \begin{center}
  \renewcommand{\arraystretch}{1}
  \begin{tabularx}{\linewidth}{|l|X|X|}
  \hline
  \textbf{Composant} & \textbf{Rôle} & \textbf{Exemple} \\[10pt]
  \hline
  Récepteur & & \\[28pt]
  \hline
  Interrupteur & & \\[28pt]
  \hline
  Fils conducteurs & & \\[28pt]
  \hline
  \end{tabularx}
  \end{center}

  \item \textit{(Document 2)} --- Explique la différence entre un circuit \textbf{ouvert} et un circuit \textbf{fermé}. Que se passe-t-il pour la lampe dans chaque cas ?

  \cadreblanc{28}

  %--- Document 3 ---
  \item \textit{(Document 3)} --- Cite \textbf{deux dangers} liés à l'électricité et explique brièvement chacun d'eux.

  \cadreblanc{30}

  \item \textit{(Document 3)} --- Cite \textbf{trois règles de sécurité} à respecter avec l'électricité.

  \cadreblanc{28}

  %--- Synthèse ---
  \item \textit{(Documents 1 et 3)} --- Le réseau électrique domestique est à 230\,V. Pourquoi est-il \textbf{interdit} de toucher une prise électrique avec les mains mouillées ? Que risque-t-on ?

  \cadreblanc{28}

  \item \textit{(Documents 2 et 3)} --- Un élève réalise un circuit avec une pile de 4,5\,V et une lampe. Est-ce que ce montage présente un danger pour lui ? Justifie ta réponse en t'appuyant sur les documents.

  \cadreblanc{28}

\end{enumerate}

\end{document}
